\documentclass[11pt]{article}
\usepackage{ap-psych}
\title{Research Methods, Statistics \& the FRQs}
\author{Paul Wang}
\date{November 2026}
\docpart{Part 0 \textperiodcentered\ Study Guide}

\begin{document}
\maketitle

\begin{abstract}
The 2024 redesign removed research methods as a stand-alone unit and spread it through all five units: stimulus-based MCQs and both free-response questions (the Article Analysis Question and the Evidence-Based Question) are built on studies. This guide covers the exam format, the research methods, sampling and assignment, threats to validity, descriptive and inferential statistics, ethics, and a part-by-part playbook for the AAQ and EBQ. Each section ends with check questions; the answer key is on the last page.
\end{abstract}

\scoringline{Check questions: \blank\ / \qtotal \qquad Target: 12+ before starting Unit 1 \qquad Each miss \(\rightarrow\) reread the section named in the Rule column}

\section{The exam at a glance}

\begin{compare}{Section}{Format}{What it rewards}
I. Multiple choice & 75 questions, 4 choices, 90 min, 66.7\% of score. Stand-alone items plus sets built on a shared study or data display. & Applying a concept to a new scenario; reading graphs and tables; spotting the method, variable or flaw. \\
II. AAQ (FRQ 1) & One summarized peer-reviewed study; six parts (A--F); 7 points; 25 min incl. 10-min reading. & Method, operational definition, statistic, ethics, generalizability, argument. \\
II. EBQ (FRQ 2) & Three summarized studies on one topic; claim + evidence + reasoning; 7 points; 45 min incl. 15-min reading. & A defensible claim backed by two cited sources and explained with course concepts. \\
\end{compare}

All five units are weighted 15--25\% of the MCQ section, so no unit can be skipped. Pacing is 72 seconds per MCQ on average; stimulus sets take longer, so bank time on the one-line definitional items.

\section{Research methods}

\begin{terms}
Experiment & The researcher \emph{manipulates} an independent variable (IV), \emph{randomly assigns} participants to conditions, and measures a dependent variable (DV). Only design that supports cause-and-effect. \\
Correlational study & Measures two or more variables as they naturally occur. Shows direction and strength of a relationship, never causation. \\
Case study & In-depth study of one person or a small group (often rare cases such as brain injury). Rich detail; cannot generalize. \\
Naturalistic observation & Records behavior in its natural setting without intervening. High realism; no control, observer effects possible. \\
Survey / self-report & Questionnaires or interviews. Fast and large; vulnerable to wording effects and social desirability bias. \\
Meta-analysis & Statistically combines results of many studies on one question to estimate an overall effect. \\
Longitudinal & Follows the \emph{same} people over time. Shows change within people; slow, attrition. \\
Cross-sectional & Compares \emph{different} age groups at one time. Fast; confounded by cohort (generation) effects. \\
Quasi-experiment & Compares groups that were not randomly assigned (e.g., smokers vs. nonsmokers). No causal claim. \\
\end{terms}

\begin{trap}
A study is an experiment only if the researcher assigns the IV. ``Researchers compared students who chose to study with music to those who did not'' is correlational or quasi-experimental, even though there are two groups and a DV.
\end{trap}

\subsection{Variables and hypotheses}
A \textbf{hypothesis} is a testable prediction derived from a theory. The \textbf{operational definition} states exactly how a variable is manipulated or measured (``anxiety = score on the GAD-7''; ``sleep deprivation = 4 hours of sleep the night before''). Good operational definitions make \textbf{replication} possible.

The \textbf{experimental group} receives the treatment; the \textbf{control group} does not (or receives a placebo). A \textbf{confounding variable} is any factor other than the IV that differs between groups and could explain the DV.

\begin{q}{3}{RM-2 Methods}{Researcher assigned sleep amount (IV) at random; only an experiment allows a causal claim.}
Participants are randomly assigned to sleep either 4 or 8 hours, then take a memory test. What type of study is this, and what is the IV?
\ch{Correlational; memory score}{Experiment; memory score}{Experiment; hours of sleep}{Quasi-experiment; hours of sleep}
\end{q}

\begin{q}{1}{RM-2 Methods}{Different ages measured once \(\rightarrow\) cross-sectional; cohort differences confound age.}
A researcher tests the vocabulary of 20-, 40-, and 60-year-olds in a single week. A critic says differences may reflect the different schooling each generation received. The design and its flaw are
\ch{cross-sectional; cohort effects}{longitudinal; attrition}{cross-sectional; experimenter bias}{longitudinal; practice effects}
\end{q}

\begin{q}{4}{RM-2 Variables}{An operational definition must state a concrete, measurable procedure.}
Which is the best operational definition of ``aggression'' for a playground study?
\ch{Hostile feelings toward peers}{How angry the children seem to the observer}{A tendency to dominate others}{The number of hits, kicks, or shoves recorded in a 30-minute recess}
\end{q}

\section{Sampling, assignment, and validity}

\begin{compare}{}{Random \emph{sampling}}{Random \emph{assignment}}
Purpose & Every member of the population has an equal chance of being chosen. & Every participant has an equal chance of being placed in any condition. \\
Protects & \textbf{Generalizability} (external validity). & \textbf{Causal inference} (internal validity): groups equal at the start. \\
Missing it means & Results may not apply beyond this sample. & Pre-existing differences may explain the DV. \\
\end{compare}

\begin{terms}
Population vs. sample & The group you want to describe vs. the subset you actually study. \\
Representative sample & Matches the population on key traits; a \textbf{convenience sample} (e.g., intro-psych students) usually does not. \\
Sampling bias & A sample that systematically over- or under-represents some group. WEIRD samples (Western, educated, industrialized, rich, democratic) limit generalization. \\
Placebo / placebo effect & An inert treatment; improvement caused by expectations alone. \\
Single-blind & Participants do not know their condition. \\
Double-blind & Neither participants nor the researchers who interact with them know conditions. Controls placebo effects \emph{and} experimenter bias. \\
Experimenter bias & Researcher expectations unintentionally influence results. \\
Social desirability bias & Answering in ways that look good rather than truthfully. \\
Hawthorne effect & Behavior changes because people know they are being observed. \\
Replication & Repeating a study (ideally with new participants and settings) to see if results hold. \\
\end{terms}

\begin{q}{2}{RM-3 Assignment}{Random assignment equalizes groups so differences can be attributed to the IV.}
Which practice best protects a study's ability to conclude that the IV caused the change in the DV?
\ch{Random sampling from the population}{Random assignment to conditions}{A large sample size}{Using a survey instead of an observation}
\end{q}

\begin{q}{3}{RM-3 Validity}{Double-blind removes both participant expectations and experimenter bias.}
In a drug trial, neither the patients nor the nurses who rate symptoms know who received the drug. This procedure mainly controls for
\ch{sampling bias and attrition}{cohort effects}{placebo effects and experimenter bias}{the Hawthorne effect only}
\end{q}

\begin{q}{1}{RM-3 Generalizability}{A convenience sample of one group limits generalization to the broader population.}
A study of stress coping used 80 volunteers from one university's psychology course. The strongest critique concerns
\ch{generalizability to the broader population}{whether there was an operational definition}{whether the study was double-blind}{the use of a meta-analysis}
\end{q}

\section{Statistics}

\subsection{Descriptive statistics}
\textbf{Mean} (average), \textbf{median} (middle score), \textbf{mode} (most frequent). Extreme scores (outliers) pull the mean toward them; the median resists outliers.
\textbf{Range} = highest -- lowest. \textbf{Standard deviation (SD)} = average distance of scores from the mean; a larger SD means more spread.

\textbf{Skew}: in a \emph{positive} (right) skew, the tail points to high scores and mean \(>\) median (e.g., income). In a \emph{negative} (left) skew, the tail points to low scores and mean \(<\) median (e.g., an easy test).

\textbf{Normal curve}: symmetrical; mean = median = mode. About 68\% of scores fall within \(\pm1\) SD, 95\% within \(\pm2\) SD, 99.7\% within \(\pm3\) SD. A \textbf{percentile rank} gives the percent of scores at or below a value (mean = 50th percentile; +1 SD \(\approx\) 84th).

\textbf{Correlation coefficient (\(r\))}: from \(-1.00\) to \(+1.00\). The sign gives direction; the absolute value gives strength. \(r = -.80\) is stronger than \(r = +.50\). A \textbf{scatterplot} shows each person as a point. An \textbf{illusory correlation} is a perceived relationship that does not exist.

\begin{trap}
Correlation does not imply causation: a third variable may drive both (ice-cream sales and drownings both rise with summer heat), or the direction may run the other way.
\end{trap}

\subsection{Inferential statistics}
\textbf{Statistical significance} asks whether a result is unlikely to be due to chance; psychology's usual cutoff is \(p < .05\) (less than a 5\% chance the difference arose by chance alone). \textbf{Effect size} measures how \emph{large} the difference or relationship is. A result can be significant but small, especially with a huge sample.

\textbf{Quantitative} data are numerical; \textbf{qualitative} data are descriptive (interview themes). Likert scales (1--5 agreement) turn attitudes into numbers.

\begin{q}{4}{RM-4 Correlation}{Strength = absolute value; |−.72| is the largest.}
Which correlation coefficient indicates the strongest relationship?
\ch{\(+.45\)}{\(+.08\)}{\(-.30\)}{\(-.72\)}
\end{q}

\begin{q}{2}{RM-4 Skew}{One extreme high score drags the mean up; the median is resistant, so mean > median (positive skew).}
Nine employees earn \$40{,}000--\$60{,}000; the CEO earns \$2 million. Which statement is accurate?
\ch{The mode is pulled toward the CEO's salary.}{The mean is higher than the median; the distribution is positively skewed.}{The median is higher than the mean; the distribution is negatively skewed.}{The mean and median are equal because the sample is small.}
\end{q}

\begin{q}{3}{RM-4 Normal curve}{+1 SD on a normal curve \(\approx\) 84th percentile (50 + 34).}
On a normally distributed test with mean 100 and SD 15, a score of 115 is at roughly the
\ch{50th percentile}{68th percentile}{84th percentile}{98th percentile}
\end{q}

\begin{q}{1}{RM-4 Significance}{p < .05 means the result is unlikely due to chance; it says nothing about size or importance.}
A study reports \(p = .01\). The most accurate interpretation is that
\ch{the difference between groups is unlikely to have occurred by chance alone}{the IV had a large effect on the DV}{1\% of participants showed the effect}{the study has been replicated}
\end{q}

\section{Ethics}

Research is reviewed by an \textbf{Institutional Review Board (IRB)} following APA guidelines:
\begin{itemize}
\item \textbf{Informed consent} (and \textbf{assent} from minors, with parental consent).
\item \textbf{Protection from harm}; risks minimized and justified by benefits.
\item \textbf{Confidentiality} (data not linked to identities) or anonymity.
\item \textbf{Right to withdraw} at any time without penalty.
\item \textbf{Deception} only when necessary and not harmful, followed by \textbf{debriefing}: explaining the true purpose afterward.
\end{itemize}
Animal research must also follow care guidelines and be justified.

\begin{q}{2}{RM-5 Ethics}{Explaining the hidden purpose afterward is debriefing.}
After a study that misled participants about its purpose, the researchers explain the true hypothesis and why deception was needed. This is
\ch{informed consent}{debriefing}{confidentiality}{the right to withdraw}
\end{q}

\begin{q}{4}{RM-5 Ethics}{Reporting only group averages with codes protects participant identities (confidentiality).}
Participants' names are replaced by ID codes and only group averages are published. Which guideline is applied?
\ch{Informed consent}{Debriefing}{Protection from harm}{Confidentiality}
\end{q}

\section{FRQ playbook}

\subsection{Article Analysis Question (AAQ), 7 points}
\begin{compare}{Part}{It asks you to}{Full-credit move}
A. Research method & Identify the method used. & Name it precisely (``experiment,'' ``correlational,'' ``meta-analysis''); no explanation needed. \\
B. Variable & State the operational definition of a named variable. & Copy the measurement procedure from the source with its numbers. \\
C. Statistic & Describe what a reported statistic or data point means. & Translate it into plain language tied to the study's variables. \\
D. Ethics & Identify how an ethical guideline was applied. & Name the guideline \emph{and} point to what the researchers did. \\
E. Generalizability & Explain the extent to which findings generalize. & Use a specific feature of the sample (size, age, country, volunteers). \\
F. Argumentation (2 pts) & Explain how the findings support or refute a claim, using a course concept. & State support/refute, then cite a specific result \emph{and} link it to the concept. \\
\end{compare}

\subsection{Evidence-Based Question (EBQ), 7 points}
You read three summarized studies on one topic and write an argument.
\begin{enumerate}[label=\textbf{\arabic*.}]
\item \textbf{Claim.} Specific and defensible (``Regular aerobic exercise reduces symptoms of depression in adults''), not a restatement of the prompt.
\item \textbf{Evidence.} Two pieces from \emph{two different sources}, cited by author and year, each with a concrete finding.
\item \textbf{Reasoning.} For each piece, explain \emph{how} it supports the claim and connect it to an AP Psychology concept (e.g., neurotransmitters, operant conditioning, the biopsychosocial model).
\end{enumerate}

\begin{keyidea}[Write-it-right checklist]
Use the source's words for variables and numbers. Answer each part separately and label it (A, B, C\dots). Do not define a concept in the abstract: apply it to the study. Never contradict yourself; a correct answer followed by a wrong one can lose the point.
\end{keyidea}

\begin{q}{3}{RM-6 AAQ}{Part E rewards a specific sample feature linked to who the results may not apply to.}
For AAQ Part E (generalizability), which response is most likely to earn the point?
\ch{``The study is not generalizable because it is only one study.''}{``More research is needed.''}{``Because all 64 participants were college students in Japan, the results may not apply to older adults or other cultures.''}{``The results are generalizable because the study was an experiment.''}
\end{q}

\begin{q}{1}{RM-6 EBQ}{A claim must be specific and arguable from the sources, not a topic label or a question.}
Which is the strongest EBQ claim for sources on social media use and adolescent sleep?
\ch{``Heavy nighttime social media use shortens adolescents' sleep duration.''}{``Social media and sleep are both important topics.''}{``Is social media good or bad?''}{``Some studies looked at social media.''}
\end{q}

\makekey
\end{document}
